% chktex-file 3
% chktex-file 8
\section{Stress-energy audit and energy conditions}\label{sec:stress-energy-audit}

\subsection{Einstein tensor of the longitudinal shift}

The effective stress-energy tensor requested for the four-dimensional slice is
\begin{equation}
  T_{\mu\nu}^{\mathrm{eff}}
  =
  \frac{1}{8\pi G_N}
  \left(
    G_{\mu\nu}[g]
    +
    \Lambda_{\mathrm{holo}}g_{\mu\nu}
  \right),
  \qquad
  \Lambda_{\mathrm{holo}}
  =
  1.08913883\times10^{-52}\;\mathrm{m^{-2}}.
  \label{eq:effective-stress-energy-definition}
\end{equation}
Energy-condition claims must be checked using the Lorentzian metric, not the
positive Euclidean Gram representative.  In the dimensionless temporal
coordinate \(x^0=ct\), write
\begin{equation}
  b(x)
  =
  \frac{\beta_x(x)}{c},
  \qquad
  g_{\mu\nu}
  =
  \begin{pmatrix}
    -1+b^2 & b & 0 & 0 \\
    b & 1 & 0 & 0 \\
    0 & 0 & 1 & 0 \\
    0 & 0 & 0 & 1
  \end{pmatrix}.
  \label{eq:section5-shift-metric}
\end{equation}
The inverse and determinant are
\begin{equation}
  g^{\mu\nu}
  =
  \begin{pmatrix}
    -1 & b & 0 & 0 \\
    b & 1-b^2 & 0 & 0 \\
    0 & 0 & 1 & 0 \\
    0 & 0 & 0 & 1
  \end{pmatrix},
  \qquad
  \det g
  =
  -1.
  \label{eq:section5-inverse-and-determinant}
\end{equation}
Primes denote derivatives with respect to \(x\).  A direct Levi-Civita
calculation gives the nonzero connection coefficients
\begin{equation}
  \begin{aligned}
  \Gamma^{0}{}_{00}
  &=
  -b^2b',
  &
  \Gamma^{0}{}_{01}
  &=
  \Gamma^{0}{}_{10}
  =
  -bb',
  &
  \Gamma^{0}{}_{11}
  &=
  -b',
  \\
  \Gamma^{1}{}_{00}
  &=
  (b^2-1)bb',
  &
  \Gamma^{1}{}_{01}
  &=
  \Gamma^{1}{}_{10}
  =
  b^2b',
  &
  \Gamma^{1}{}_{11}
  &=
  bb'.
  \end{aligned}
  \label{eq:shift-christoffel-symbols}
\end{equation}
Define
\begin{equation}
  Q(x)
  =
  b(x)b''(x)
  +
  [b'(x)]^2
  =
  \frac{1}{2}
  \frac{d^2}{dx^2}
  b^2(x).
  \label{eq:shift-curvature-functional}
\end{equation}
The Ricci tensor and scalar are
\begin{equation}
  R_{\mu\nu}
  =
  \begin{pmatrix}
    (b^2-1)Q & bQ & 0 & 0 \\
    bQ & Q & 0 & 0 \\
    0 & 0 & 0 & 0 \\
    0 & 0 & 0 & 0
  \end{pmatrix},
  \qquad
  R
  =
  2Q.
  \label{eq:shift-ricci-tensor}
\end{equation}
Consequently the Einstein tensor is particularly simple:
\begin{equation}
  G_{\mu\nu}
  =
  \begin{pmatrix}
    0 & 0 & 0 & 0 \\
    0 & 0 & 0 & 0 \\
    0 & 0 & -Q & 0 \\
    0 & 0 & 0 & -Q
  \end{pmatrix}.
  \label{eq:shift-einstein-tensor}
\end{equation}
Equation~\eqref{eq:shift-einstein-tensor} agrees with the flat-interior result:
on an exact plateau \(b'=b''=0\), so \(Q=0\) and
\(G_{\mu\nu}=0\).

\subsection{Null energy condition}

Let \(k^\mu\) be any null vector,
\begin{equation}
  g_{\mu\nu}k^\mu k^\nu
  =
  0.
  \label{eq:null-vector-condition}
\end{equation}
The cosmological term drops out of the null contraction, and
Eq.~\eqref{eq:shift-einstein-tensor} yields
\begin{equation}
  T_{\mu\nu}^{\mathrm{eff}}
  k^\mu k^\nu
  =
  -\frac{Q(x)}{8\pi G_N}
  \left[
    (k^y)^2+(k^w)^2
  \right].
  \label{eq:nec-contraction}
\end{equation}
Null vectors with nonzero transverse components exist at every regular point.
Therefore the necessary and sufficient pointwise condition for the NEC is
\begin{equation}
  T_{\mu\nu}^{\mathrm{eff}}
  k^\mu k^\nu
  \geq0
  \quad
  \text{for all null }k^\mu
  \qquad
  \Longleftrightarrow
  \qquad
  Q(x)\leq0.
  \label{eq:nec-condition-q}
\end{equation}

For a localized shift, impose the asymptotic conditions
\begin{equation}
  \lim_{x\to\pm\infty}b(x)
  =
  0,
  \qquad
  \lim_{x\to\pm\infty}b'(x)
  =
  0.
  \label{eq:localized-shift-boundary-conditions}
\end{equation}
Since \(Q=(bb')'\),
\begin{equation}
  \int_{-\infty}^{+\infty}
  Q(x)\,dx
  =
  \left[
    b(x)b'(x)
  \right]_{-\infty}^{+\infty}
  =
  0.
  \label{eq:integrated-q-zero}
\end{equation}
If the NEC held everywhere, Eq.~\eqref{eq:nec-condition-q} would give
\(Q\leq0\) everywhere.  Together with
Eq.~\eqref{eq:integrated-q-zero}, continuity would force
\begin{equation}
  Q(x)
  =
  0
  \quad
  \text{for every }x.
  \label{eq:q-identically-zero}
\end{equation}
Equation~\eqref{eq:shift-curvature-functional} would then imply that \(b^2\)
is affine in \(x\).  The two asymptotic conditions in
Eq.~\eqref{eq:localized-shift-boundary-conditions} force that affine function
to vanish:
\begin{equation}
  Q\equiv0
  \quad\Longrightarrow\quad
  b^2(x)=A x+B
  \quad\Longrightarrow\quad
  A=B=0
  \quad\Longrightarrow\quad
  b(x)\equiv0.
  \label{eq:nec-localized-no-go}
\end{equation}
Hence a nontrivial, smooth, localized shift of the form
Eq.~\eqref{eq:section5-shift-metric} cannot satisfy the NEC everywhere.
Longitudinal null vectors with \(k^y=k^w=0\) saturate the NEC, but transverse
null directions detect the violation wherever \(Q>0\).

\subsection{Weak energy condition}

Let \(u^\mu\) be a unit timelike vector,
\begin{equation}
  g_{\mu\nu}u^\mu u^\nu
  =
  -1.
  \label{eq:unit-timelike-condition}
\end{equation}
Using Eqs.~\eqref{eq:effective-stress-energy-definition} and
~\eqref{eq:shift-einstein-tensor}, the measured effective energy density is
\begin{equation}
  T_{\mu\nu}^{\mathrm{eff}}
  u^\mu u^\nu
  =
  \frac{1}{8\pi G_N}
  \left\{
    -Q(x)
    \left[
      (u^y)^2+(u^w)^2
    \right]
    -
    \Lambda_{\mathrm{holo}}
  \right\}.
  \label{eq:wec-contraction}
\end{equation}
The comoving Eulerian observer is
\begin{equation}
  n^\mu
  =
  (1,-b,0,0),
  \qquad
  g_{\mu\nu}n^\mu n^\nu
  =
  -1.
  \label{eq:eulerian-observer-section5}
\end{equation}
For this observer,
\begin{equation}
  T_{\mu\nu}^{\mathrm{eff}}
  n^\mu n^\nu
  =
  -\frac{\Lambda_{\mathrm{holo}}}
  {8\pi G_N}.
  \label{eq:eulerian-wec-density}
\end{equation}
Thus, with a positive \(\Lambda_{\mathrm{holo}}\) included on the
stress-energy side exactly as in
Eq.~\eqref{eq:effective-stress-energy-definition}, the WEC is already violated
for the comoving observer.  If the cosmological term is retained on the
geometric side and excluded from the material stress tensor, one sets
\(\Lambda_{\mathrm{holo}}=0\) in
Eq.~\eqref{eq:wec-contraction}.  The WEC for all timelike observers then again
requires
\begin{equation}
  Q(x)\leq0,
  \label{eq:wec-condition-zero-lambda}
\end{equation}
and the localized no-go theorem
Eq.~\eqref{eq:nec-localized-no-go} applies unchanged.

\subsection{Comparison with the classical Alcubierre source}

For the usual three-dimensional Alcubierre profile
\(\beta_x=-v_s f(r_s)\), the Eulerian energy density contains the explicitly
negative transverse-gradient term
\begin{equation}
  \rho_{\mathrm{Alc}}
  =
  -\frac{c^4}{32\pi G_N}
  \frac{v_s^2}{c^2}
  \left[
    (\partial_y f)^2
    +
    (\partial_w f)^2
  \right]
  \leq0.
  \label{eq:classical-alcubierre-energy-density}
\end{equation}
The restricted one-dimensional shift used in
Eq.~\eqref{eq:section5-shift-metric} removes these transverse derivatives.
For \(\Lambda_{\mathrm{holo}}=0\), its coordinate component and Eulerian
energy density are
\begin{equation}
  T_{00}^{\mathrm{eff}}
  =
  0,
  \qquad
  T_{\mu\nu}^{\mathrm{eff}}
  n^\mu n^\nu
  =
  0.
  \label{eq:one-dimensional-zero-eulerian-density}
\end{equation}
Boundary-driven longitudinal shear therefore avoids the specific negative
\(T_{00}\) term in Eq.~\eqref{eq:classical-alcubierre-energy-density} for
this reduced ansatz.  It does not, however, establish universal energy
condition compliance: the transverse pressures in
Eq.~\eqref{eq:shift-einstein-tensor} still violate the NEC and WEC somewhere
for every nontrivial localized profile.

Accordingly, a claim of complete SHBT energy-condition compliance requires
additional dynamics beyond
Eq.~\eqref{eq:effective-stress-energy-definition}.  Such a completion may be
parameterized by a boundary-response tensor
\(\mathcal B_{\mu\nu}\):
\begin{equation}
  8\pi G_N
  T_{\mu\nu}^{\mathrm{phys}}
  =
  G_{\mu\nu}
  +
  \Lambda_{\mathrm{holo}}g_{\mu\nu}
  +
  \mathcal B_{\mu\nu}.
  \label{eq:boundary-response-completion}
\end{equation}
For the NEC and WEC to hold, this tensor must satisfy the explicit
inequalities
\begin{equation}
  \mathcal B_{\mu\nu}k^\mu k^\nu
  \geq
  Q
  \left[
    (k^y)^2+(k^w)^2
  \right]
  \quad
  \text{for all null }k^\mu,
  \label{eq:boundary-response-nec-requirement}
\end{equation}
and
\begin{equation}
  \mathcal B_{\mu\nu}u^\mu u^\nu
  \geq
  Q
  \left[
    (u^y)^2+(u^w)^2
  \right]
  +
  \Lambda_{\mathrm{holo}}
  \quad
  \text{for all unit timelike }u^\mu.
  \label{eq:boundary-response-wec-requirement}
\end{equation}
Until \(\mathcal B_{\mu\nu}\) is derived independently from the boundary
theory, Eqs.~\eqref{eq:nec-localized-no-go} and
~\eqref{eq:eulerian-wec-density} are the rigorous conclusions of the
Einstein-tensor audit.